\documentclass[10pt,a4paper]{amsart}
\usepackage[margin=19mm]{geometry}
\usepackage{amsmath,amssymb,mathtools}
\usepackage[scr=rsfs,cal=euler]{mathalfa}
\usepackage{xfrac}
\usepackage[unicode,hidelinks,bookmarks=false]{hyperref}
\newcommand{\ie}{i.e.\ }
\newcommand{\cf}{cf.\ }
\newcommand{\resp}{respectively\ }
\newcommand{\Subsection}{\subsection}
\providecommand{\nobreakdash}{\nobreak\hbox{-}}
\setlength{\emergencystretch}{2em}
\multlinegap0pt
\usepackage{lmodern}
\usepackage{enumitem}
\usepackage{mathtools,mathrsfs}
\newtheorem{theorem}{Theorem}[section]
\newtheorem{corollary}[theorem]{Corollary}
\newtheorem{lemma}[theorem]{Lemma}
\newtheorem{proposition}[theorem]{Proposition}
\newtheorem*{claim}{Claim}
\newtheorem{definition}[theorem]{Definition}
\newtheorem{remark}[theorem]{Remark}
\newcommand{\T}{\mathbb T}
\newcommand{\N}{\mathbb N}
\newcommand{\Z}{\mathbb Z}
\newcommand{\Q}{\mathbb Q}
\newcommand{\cfrak}{\mathfrak c}
\newcommand{\Hom}{\operatorname{Hom}}
\newcommand{\supp}{\operatorname{supp}}
\newcommand{\Rel}{\operatorname{Rel}}
\newcommand{\htop}{\operatorname{ht}}
\begin{document}
\title[The Wallace problem and countably compact torsion-free Abeli]{The Wallace problem and countably compact torsion-free Abelian groups in ZFC}
\author{Secondary 20K20, 54A20 ęywords countably compactness, topological groups, torsion-free Abelian groups,nontrivial convergent sequence; Juliane Trianon Fraga; Vinicius de Oliveira Rodrigues}
\date{}
\maketitle
\noindent\emph{Source and licence.} The Wallace problem and countably compact torsion-free Abelian groups in ZFC. Version arXiv:2608.17317v1 (CC BY 4.0). This material is distributed under the Creative Commons Attribution 4.0 International licence, http://creativecommons.org/licenses/by/4.0/, as recorded in the arXiv OAI licence metadata for this exact version and shown on the abstract page https://arxiv.org/abs/2608.17317v1. Full rights notice: licenses/arxiv-2608.17317-CC-BY-4.0.txt.\par\smallskip
\noindent\textit{Editorial scope.} Counted unit: the original \texttt{lemma} environment \texttt{lem:fusion-step} at source lines 926--958 together with its complete proof at lines 960--1065; the delivered document keeps source lines 203--1066 so that every referenced label and every citation resolves inside it. Native editable source is the paper's own concatenated TeX; the AMS wrapper is editorial formatting only. No publisher class, upstream script or unrelated artwork is redistributed.
\par\smallskip\noindent\textit{Reference note.} This article ships no compiled bibliography (biblatex); the entries delivered here are composed from the author's own BibTeX (.bib) fields, with no field invented and no entry dropped.
\section{Preliminaries and the topological reduction}\label{sec:preliminaries}

We write $\N=\{0,1,2,\ldots\}$, regard each $m\in\N$ as the finite ordinal
$\{i:i<m\}$, and index every $m$-tuple by $i<m$.  We use additive notation.
The circle group
$\T=\mathbb R/\mathbb Z$ carries the quotient metric induced by the Euclidean
metric on $\mathbb R$, determined by
\[
 \|x+\mathbb Z\|=\min_{n\in\mathbb Z}|x-n|
 \qquad(x\in\mathbb R).
\]
If $(x_n:n\in\N)$ is a sequence in a Hausdorff topological space and $p$ is
an ultrafilter on $\N$, then
$p\text{-}\lim_n x_n=x$ means that
\[
 \text{for every neighborhood }U\text{ of }x,\ \{n:x_n\in U\}\in p.
\]
Every ultrafilter on a compact Hausdorff space converges to a unique point.
An ultrafilter on $\N$ is free when it
contains the cofinite filter.

The circle is divisible and therefore injective in the category of Abelian
groups.  Consequently a homomorphism from a subgroup of an Abelian group to
$\T$ extends to the whole group
\cite[Theorem~21.1]{Fuchs1970}.
We write $\Hom(G,\T)$ for the set of group homomorphisms from an Abelian
group $G$ to $\T$.

As usual, $\Q^{(\cfrak)}=\{z\in\Q^{\cfrak}:|\supp(z)|<\infty\}$ denotes the
direct sum of $\cfrak$ copies of $\Q$, where
$\supp(z)=\{\xi<\cfrak:z(\xi)\neq0\}$ is the support of $z$.

We adopt the convention that, if $D\subseteq\cfrak$, then $\Q^{(D)}$
denotes the subgroup of $\Q^{(\cfrak)}$ consisting of all vectors supported
in $D$.  Thus, each element of $\Q^{(D)}$ is a function from $\cfrak$ into
$\Q$ that vanishes outside $D$.

For each $\xi<\cfrak$, let $e_\xi\in\Z^{(\cfrak)}$ denote the vector
supported at $\xi$ and defined by
\[
 e_\xi(\eta)=
 \begin{cases}
  1,&\text{if $\eta=\xi$},\\
  0,&\text{if $\eta\neq\xi$}.
 \end{cases}
\]

\begin{lemma}\label{lem:torsion-free-coordinatization}
Let $G$ be a torsion-free Abelian group of cardinality $\cfrak$.  Then $G$ may
be identified with a subgroup of $\Q^{(\cfrak)}$ in such a way that
\[
 \Z^{(\cfrak)}\leq G\leq\Q^{(\cfrak)}.
\]
In particular, for every $\xi<\cfrak$, the vector $e_\xi$ belongs to $G$.
\end{lemma}

\begin{proof}
The canonical homomorphism
\[
 G\longrightarrow\Q\otimes_{\Z}G,
 \qquad x\longmapsto1\otimes x,
\]
is injective because $G$ is torsion-free.  Its image spans the rational vector
space $\Q\otimes_{\Z}G$. As $\mathbb Q$ is countable, the dimension of this vector space is $\cfrak$ because $G$ has cardinality $\cfrak$.

Choose a basis $(1\otimes b_\xi)_{\xi<\cfrak}$ from the spanning set given by
the image of $G$, and identify $\Q\otimes_{\Z}G$ with
$\Q^{(\cfrak)}$ by sending $1\otimes b_\xi$ to $e_\xi$.  Under the resulting
embedding of $G$, each $b_\xi$ is identified with $e_\xi$, and hence
$\Z^{(\cfrak)}\leq G\leq\Q^{(\cfrak)}$.
\end{proof}

For a finite tuple $z=(z_0,\ldots,z_{m-1})$ in an Abelian group put
\[
 \Rel(z)=\left\{a\in\Z^m:\sum_{i<m} a_i z_i=0\right\},
 \qquad
 \htop(a)=\max\bigl(\{0\}\cup\{|a_i|:i<m\}\bigr).
\]

In a $T_1$ space, countable
compactness is equivalent to the assertion that every countably infinite
subset has an accumulation point
\cite[Theorem~3.10.3(v)]{Engelking1989}.  Also, if $u:\N\to X$ is injective,
$p$ is free, and $p\text{-}\lim_n u(n)=x$, then $x$ is an accumulation point
of $u[\N]$.

The following observation will be applied to show that the group topologies
we construct are countably compact.
This is a standard result and we include a proof for completeness.

\begin{lemma}\label{lem:topological-reduction}
Let $G$ be a Hausdorff topological group.  Suppose that for every injective
$s:\N\to G$ there exist a strictly increasing map $\varphi:\N\to\N$, a free
ultrafilter $p$ on $\N$, and a point $b\neq0$ such that
\begin{equation}\label{eq:nonzero-limit-property}
 p\text{-}\lim_n s(\varphi(n))=b.
\end{equation}
Then $G$ is countably compact and every convergent sequence in $G$ is
eventually constant.
\end{lemma}

\begin{proof}
Given a countably infinite set $A\subseteq G$, enumerate it injectively by
$s$ and apply the hypothesis.  It follows that $A$ has an accumulation point,
so $G$ is countably compact.

Now suppose that an injective sequence $s$ converges to $x$.  Define
$t:\N\to G$ by $t(n)=s(n)-x$.
The sequence $t$ is injective and converges to $0$.  Apply the hypothesis
to $t$, obtaining $\varphi,p$ and $b\neq0$.  Since $\varphi$ is strictly
increasing, $t\circ\varphi$ still converges to $0$, and since $p$ is
free, it follows that $p\text{-}\lim_n t(\varphi(n))=0$.  This contradicts
uniqueness of limits in a Hausdorff space, so $s$ cannot be injective.

Finally let $s$ be any sequence converging to $x$.  If its range is finite,
then
$s(n)=x$ eventually.  If its range is infinite, choose recursively a
strictly increasing subsequence with pairwise distinct values.  That
subsequence still converges to $x$, contradicting the preceding paragraph.
\end{proof}

Suppose now that $(\psi_j)_{j\in J}$ is a family of homomorphisms in
$\Hom(G,\T)$ that separates points: whenever $x\neq y$, there exists $j\in J$
such that $\psi_j(x)\neq\psi_j(y)$.  Define
\[
 \Delta:G\longrightarrow\T^J,
 \qquad \Delta(x)=(\psi_j(x))_{j\in J}.
\]
Give $G$ the coarsest topology that makes every $\psi_j$ continuous, that is,
the initial topology induced by the family $(\psi_j)_{j\in J}$. This is a
Hausdorff group topology, and $\Delta$ is a topological embedding
onto its image.  Thus, once these homomorphisms have been constructed, only
\eqref{eq:nonzero-limit-property} remains to obtain all the topological
conclusions.


\section{Uniform finite approximation by circle-valued homomorphisms}

A basic ingredient in the construction is the
following finite interpolation problem.  Given elements
$z_0,\ldots,z_{m-1}$ of an Abelian group and prescribed target values
$t_0,\ldots,t_{m-1}\in\T$, we seek a single homomorphism from the ambient
group to $\T$ whose value at each $z_i$ is ``close'' to $t_i$.

The following Kronecker-type lemma shows
that, for approximate realization to a fixed accuracy, it is enough to check
compatibility only with the relations whose coefficients are bounded by a
suitable finite constant.

\begin{lemma}[Uniform finite Kronecker lemma]\label{lem:uniform-kronecker}
For every integer $m\geq1$ and every $\varepsilon>0$ there exists $q\in\N$
with the following property.

For every Abelian group $G$, for every
$z=(z_i)_{i<m}\in G^m$, and for every
$t=(t_i)_{i<m}\in\T^m$, if for every relation
$a=(a_i)_{i<m}\in\Rel(z)$ with $\htop(a)\leq q$ we have
$\sum_{i<m}a_i t_i=0$, then there exists $\psi\in\Hom(G,\T)$ such that
\[
 \max_{i<m}\|\psi(z_i)-t_i\|<\varepsilon.
\]
\end{lemma}
We note that the integer $q$ is independent of $G,z$, and $t$.

We dedicate this section to proving Lemma~\ref{lem:uniform-kronecker}.
First, we isolate the analytic ingredients that will be needed.  For each
integer $m\geq1$, on $\T^m$ we use the maximum metric
\[
 \rho_\infty(u,v)=\max_{i<m}\|u_i-v_i\|.
\]
Let $\mathbb S^1=\{w\in\mathbb C:|w|=1\}$.  The homomorphism
$r\mapsto e^{2\pi i r}$ from $\mathbb R$ onto $\mathbb S^1$ has kernel
$\mathbb Z$, so it induces the continuous group isomorphism
\[
 E:\T=\mathbb R/\mathbb Z\longrightarrow\mathbb S^1,
 \qquad E(r+\mathbb Z)=e^{2\pi i r}.
\]
For $a=(a_i)_{i<m}\in\Z^m$, define the continuous homomorphism
\[
 \phi_a:\T^m\longrightarrow\mathbb S^1,
 \qquad \phi_a(u)=E\left(\sum_{i<m}a_i u_i\right).
\]

The only compactness result for families of functions that we need is the
following compact-metric form of Arzel\`a--Ascoli.  For a proof, see, e.g.,
\cite[Theorem~8.2.10]{Engelking1989}.

\begin{lemma}[Arzel\`a--Ascoli]\label{lem:arzela-ascoli-uniform}
Let $X$ be a compact metric space.  If
$\mathcal F\subseteq C(X,\mathbb R)$ is equicontinuous and uniformly bounded,
then its closure in the uniform norm is compact.
\end{lemma}

For the following standard complex form of the Stone--Weierstrass theorem,
see \cite[Theorem~7.33]{Rudin1976}.  Recall that a subset
$\mathcal A\subseteq C(X,\mathbb C)$ is uniformly dense in
$C(X,\mathbb C)$ if it is dense in the topology induced by the uniform norm
$\|\cdot\|_\infty$.

\begin{theorem}[Complex Stone--Weierstrass theorem]
\label{thm:complex-stone-weierstrass}
Let $X$ be a compact Hausdorff space, and let
$\mathcal A\subseteq C(X,\mathbb C)$ be a complex subalgebra.  If
$\mathcal A$ contains the constant functions, separates the points of $X$,
and is closed under complex conjugation, then $\mathcal A$ is uniformly dense
in $C(X,\mathbb C)$.
\end{theorem}

\begin{lemma}
\label{lem:trigonometric-stone-weierstrass}
For every integer $m\geq1$, the complex linear span of
$\{\phi_a:a\in\Z^m\}$ is uniformly dense in $C(\T^m,\mathbb C)$.
\end{lemma}

\begin{proof}
Their complex linear span is a unital algebra because
$\phi_0=1$ and $\phi_a\phi_b=\phi_{a+b}$, and it is closed under complex
conjugation because $\overline{\phi_a}=\phi_{-a}$.  It also separates
points: if $u\neq v$, then $u_j\neq v_j$ for some $j<m$, and the homomorphism
$\phi_{e_j}$, where $e_j$ is the $j$th standard basis vector, distinguishes
$u$ and $v$.
Theorem~\ref{thm:complex-stone-weierstrass} now applies.
\end{proof}

We will call each element of the span of $\{\phi_a:a\in\Z^m\}$ a
\emph{trigonometric polynomial} on $\T^m$.

We now combine these results to obtain the lemma below.
In its statement, $1$-Lipschitz means Lipschitz with constant $1$ with respect
to the metric $\rho_\infty$ on $\T^m$ and the usual metric on $\mathbb R$.

\begin{lemma}\label{lem:uniform-frequency-set}
For every integer $m\geq1$ and every $\eta>0$, there exists a finite set
$S\subseteq\Z^m$ such that for every $1$-Lipschitz function
$f:\T^m\to[0,1/2]$ there exist complex coefficients $(c_a)_{a\in S}$
satisfying
\[
 \left\|f-\sum_{a\in S}c_a\phi_a\right\|_\infty<\eta.
\]
\end{lemma}

\begin{proof}
Let $\mathcal L_m$ be the set of all such functions, regarded as a subset of
$C(\T^m,\mathbb R)$.  This family is uniformly bounded, equicontinuous (as
all its members are $1$-Lipschitz), and closed under uniform limits.  Hence it
is closed in $C(\T^m,\mathbb R)$ and compact by
Lemma~\ref{lem:arzela-ascoli-uniform}.

For each $f\in\mathcal L_m$, use
Lemma~\ref{lem:trigonometric-stone-weierstrass} to choose a trigonometric
polynomial $P_f$ such that $\|f-P_f\|_\infty<\eta/2$.  The uniform balls of
radius $\eta/2$ centered at the members of $\mathcal L_m$ cover
$\mathcal L_m$; choose a finite subcover with centers
$f_0,\ldots,f_N$.
For each $j\leq N$, choose a finite set $S_j\subseteq\Z^m$ such that
$P_{f_j}\in\operatorname{span}\{\phi_a:a\in S_j\}$, and set
$S=\bigcup_{j\leq N} S_j$.

If $f$ belongs to the ball centered at $f_j$, then
\[
 \|f-P_{f_j}\|_\infty
 \leq \|f-f_j\|_\infty+\|f_j-P_{f_j}\|_\infty<\eta.
\]
This proves
the assertion.
\end{proof}

For an integer $m\geq1$ and a subgroup $R\leq\Z^m$, put
\[
 H_R=\left\{u\in\T^m:\sum_{i<m}a_i u_i=0
                    \text{ for every }a\in R\right\}.
\]

Recall that if $H$ is a compact Hausdorff Abelian topological group, its
\emph{normalized Haar measure} is the unique regular Borel measure $\mu$
on $H$ such that $\mu(H)=1$ and
$\mu(x+B)=\mu(B)$ for every $x\in H$ and every Borel set $B\subseteq H$.
Thus $\mu$ is a probability measure invariant under translations; in
particular, for every continuous $g:H\to\mathbb C$ and every $x\in H$,
\[
 \int_H g(x+h)\,d\mu(h)=\int_H g(h)\,d\mu(h).
\]
Existence and uniqueness are standard; see
\cite[Chapter~VII, \S\S~2--3]{Katznelson2004}.

\begin{lemma}
\label{lem:annihilator-haar-average}
For every integer $m\geq1$ and every subgroup $R\leq\Z^m$, the set $H_R$ is
a compact subgroup of $\T^m$, and
\begin{equation}\label{eq:double-annihilator}
 \{a\in\Z^m:\phi_a(h)=1\text{ for every }h\in H_R\}=R.
\end{equation}
If $\mu_R$ is normalized Haar measure on $H_R$, then, for every
$a\in\Z^m$ and $s\in\T^m$,
\begin{equation}\label{eq:haar-homomorphism-average}
 \int_{H_R}\phi_a(s+h)\,d\mu_R(h)
 =\begin{cases}
   \phi_a(s),&a\in R,\\
   0,&a\notin R.
  \end{cases}
\end{equation}
\end{lemma}

\begin{proof}
The set $H_R$ is an intersection of kernels of continuous homomorphisms
from $\T^m$ to $\T$, so it is a closed subgroup of the compact group
$\T^m$.  The inclusion from right to left in
\eqref{eq:double-annihilator} follows from the definition of $H_R$.

For the reverse inclusion, let $\pi:\Z^m\to\Z^m/R$ be the quotient map and
suppose that $a\notin R$.  Since $\pi(a)\neq0$, there is a homomorphism
$\phi:\Z^m/R\to\T$ such that $\phi(\pi(a))\neq0$: first define it on
$\langle\pi(a)\rangle$ and then use the injectivity of $\T$ to extend it.
For $i<m$, put $u_i=\phi(\pi(e_i))$.  If $r=(r_i)_{i<m}\in R$, then
\[
 \sum_{i<m}r_i u_i
 =\phi\left(\pi\left(\sum_{i<m}r_i e_i\right)\right)
 =\phi(\pi(r))=0,
\]
so $u=(u_i)_{i<m}\in H_R$.  On the other hand,
\[
 \sum_{i<m}a_i u_i=\phi(\pi(a))\neq0.
\]
Thus $\phi_a(u)\neq1$, proving \eqref{eq:double-annihilator}.

Put
\[
 I_a=\int_{H_R}\phi_a(h)\,d\mu_R(h).
\]
If $a\in R$, then $\phi_a$ is identically $1$ on $H_R$, so $I_a=1$.  If
$a\notin R$, identity \eqref{eq:double-annihilator} gives
$h_0\in H_R$ with $\phi_a(h_0)\neq1$.  Translation invariance gives
\[
 I_a=\int_{H_R}\phi_a(h_0+h)\,d\mu_R(h)=\int_{H_R}\phi_a(h_0)\phi_a(h)\,d\mu_R(h)=\phi_a(h_0)I_a,
\]
and hence $I_a=0$.  Finally,
$\phi_a(s+h)=\phi_a(s)\phi_a(h)$, so
\[
  \int_{H_R}\phi_a(s+h)\,d\mu_R(h)=\phi_a(s)I_a,
\]
which proves \eqref{eq:haar-homomorphism-average}.
\end{proof}

Now we are ready to prove Lemma~\ref{lem:uniform-kronecker}.

\begin{proof}[Proof of Lemma~\ref{lem:uniform-kronecker}]
Set $\eta=\varepsilon/4$, and let $S\subseteq\Z^m$ be supplied by
Lemma~\ref{lem:uniform-frequency-set}.
Put
\[
 q=\max\bigl(\{1\}\cup\{\htop(a):a\in S\}\bigr).
\]

\begin{claim}
For every subgroup $R\leq\Z^m$ and every
$t\in\T^m$, if
\[
 \sum_{i<m}a_i t_i=0
 \qquad\text{for every }a\in R\cap[-q,q]^m,
\]
then $\rho_\infty(t,H_R)<\varepsilon$.
\end{claim}

\begin{proof}[Proof of the claim]
Fix a subgroup $R\leq\Z^m$ and $t\in\T^m$ such that
\begin{equation}\label{eq:short-annihilation}
 \sum_{i<m}a_i t_i=0\text{ for every }a\in R\cap[-q,q]^m.
\end{equation}
Define $d:\T^m\to\mathbb R$ by
\[
 d(u)=\rho_\infty(u,H_R)=\inf_{h\in H_R}\rho_\infty(u,h).
\]
The function $d$ is $1$-Lipschitz.
Moreover, $0\in H_R$, and hence
\[
 0\leq d(u)\leq\rho_\infty(u,0)\leq\frac12
 \qquad(u\in\T^m).
\]
Lemma~\ref{lem:uniform-frequency-set}, applied to $d$, gives coefficients
$(c_a)_{a\in S}\in\mathbb C^S$ such that the trigonometric polynomial
\[
 P(u):=\sum_{a\in S}c_a\phi_a(u)
\]
satisfies $\|P-d\|_\infty<\eta$.

For $s\in\T^m$, set
\[
 A(s)=\int_{H_R}P(s+h)\,d\mu_R(h).
\]
Equation~\eqref{eq:haar-homomorphism-average} yields
\begin{equation}\label{eq:average-polynomial}
 A(s)=\sum_{a\in S\cap R}c_{a}\phi_a(s).
\end{equation}
By the definition of $q$,
$S\cap R\subseteq R\cap[-q,q]^m$.  Thus
\eqref{eq:short-annihilation} implies $\phi_a(t)=1$ for every
$a\in S\cap R$, and \eqref{eq:average-polynomial} gives
\begin{equation}\label{eq:equal-averages}
 A(t)=\sum_{a\in S\cap R}c_a\phi_a(t)
     =\sum_{a\in S\cap R}c_a
     =\sum_{a\in S\cap R}c_a\phi_a(0)=A(0).
\end{equation}

Because $d$ vanishes on $H_R$ and $\mu_R$ is a probability measure,
\begin{equation}\label{eq:average-at-zero}
 |A(0)|
 =\left|\int_{H_R}\bigl(P(h)-d(h)\bigr)\,d\mu_R(h)\right|\leq\|P-d\|_\infty<\eta.
\end{equation}
Since $H_R$ is a subgroup and $\rho_\infty$ is translation invariant,
$d(t+h)=d(t)$ for every $h\in H_R$.  Consequently,
$d(t)=\int_{H_R}d(t+h)\,d\mu_R(h)$, and therefore
\begin{equation}\label{eq:average-at-t}
 |d(t)-A(t)|
 =\left|\int_{H_R}\bigl(d(t+h)-P(t+h)\bigr)\,d\mu_R(h)\right|
 <\eta.
\end{equation}
Combining \eqref{eq:equal-averages}--\eqref{eq:average-at-t}, we obtain
\[
 d(t)=|d(t)|
 \leq |d(t)-A(t)|+|A(0)|
 <2\eta<\varepsilon.
\]
\end{proof}

Now let $G$ be an Abelian group, let
$z=(z_i)_{i<m}\in G^m$ and
$t=(t_i)_{i<m}\in\T^m$ satisfy the hypothesis of the lemma, and take
$R=\Rel(z)$.
Then $t$ satisfies the hypothesis of the claim for this $R$.  Hence $\rho_\infty(t,H_R)<\varepsilon$, so there
exists $u\in H_R$ such that
$\rho_\infty(u,t)<\varepsilon$.  Define a map on the subgroup
$L=\langle z_i:i<m\rangle\leq G$ by
\[
 \varphi\left(\sum_{i<m}n_i z_i\right)=\sum_{i<m}n_i u_i.
\]
To see that this is well defined, suppose that $n,n'\in\Z^m$ satisfy
$\sum_{i<m}n_i z_i=\sum_{i<m}n'_i z_i$.  Then $n-n'\in\Rel(z)=R$, and since
$u\in H_R$,
\[
 \sum_{i<m}n_i u_i-\sum_{i<m}n'_i u_i
 =\sum_{i<m}(n_i-n'_i)u_i=0.
\]
Clearly, $\varphi$ is a homomorphism from $L$ to $\T$.
By injectivity of $\T$, the map $\varphi$ extends to some
$\psi\in\Hom(G,\T)$.  For each $i<m$ we have $\psi(z_i)=u_i$, and therefore
\[
 \max_{i<m}\|\psi(z_i)-t_i\|=\rho_\infty(u,t)<\varepsilon.
\]
\end{proof}

\section{Bounded integer relations and deletion}

The goal of this section is Lemma~\ref{lem:fusion-step}.  Given a
homomorphism $\psi:G\to\T$ and finite sets $A,X\subseteq G$, that lemma
produces a large subset $Y\subseteq X$ and a homomorphism $\psi':G\to\T$
which approximately preserves the values of $\psi$ on $A$ and is small on
$Y$.  Lemma~\ref{lem:uniform-kronecker} reduces this construction to
eliminating bounded-height integer relations that mix elements of $A$ with
elements of $Y$.  We now develop the finite combinatorial deletion argument
needed to achieve this.

\begin{definition}
Let $G$ be an Abelian group.  For a finite set $X\subseteq G$, put
\[
 \Rel_G(X)=
 \left\{c=(c_x)_{x\in X}\in\Z^X:
 \sum_{x\in X}c_x\,x=0\right\}.
\]
This is the set-indexed version of the notation $\Rel(z)$ introduced in
Section~\ref{sec:preliminaries}.
For $c\in\Z^X$, put
\[
 \htop(c)=\max\bigl(\{0\}\cup\{|c_x|:x\in X\}\bigr).
\]

For $M\in\N$, the set $X$ is \emph{$M$-independent} if the zero vector is
the only $c\in\Rel_G(X)$ satisfying $\htop(c)\leq M$.

For finite sets $A,Y\subseteq G$, put
\[
 \Rel_G(A,Y)=
 \left\{(b,c)\in\Z^A\times\Z^Y:
 \sum_{a\in A}b_a\,a+\sum_{y\in Y}c_y\,y=0\right\}.
\]
An element $(b,c)\in\Rel_G(A,Y)$ is \emph{mixed} if $b\neq0$ and
$c\neq0$.  We set $\htop(b,c)=\max\{\htop(b),\htop(c)\}$.
\end{definition}

\begin{lemma}\label{lem:finite-extraction}
Let $G$ be a torsion-free Abelian group and $S\subseteq G$ be infinite.
Given $M,N\in\N$,
there exists an $M$-independent subset of $S$ of cardinality $N$.
\end{lemma}

\begin{proof}
We construct, by induction on $k\leq N$, an $M$-independent set
$B_k\subseteq S$ such that $|B_k|=k$.

For $k=0$, take $B_0=\varnothing$, which is trivially
$M$-independent.  Suppose that $k<N$ and that
\[
B_k=\{x_i:i<k\}\subseteq S
\]
has already been chosen and is $M$-independent. We shall choose
$x\in S\setminus B_k$ such that $B_k\cup\{x\}$ remains
$M$-independent.

Indeed, given $x\in S\setminus B_k$, if $B_k\cup\{x\}$ is not
$M$-independent, then there exist integers $q$ and $(c_i)_{i<k}$, not all
zero, with
\[
|q|\leq M,\qquad |c_i|\leq M\quad(i<k),
\]
such that
\[
qx+\sum_{i<k} c_i x_i=0.
\]
We claim that $q\neq0$. Indeed, if $q=0$, then the family
$c\in\Z^{B_k}$ defined by $c_{x_i}=c_i$ for $i<k$ would be a nonzero
element of $\Rel_G(B_k)$ with $\htop(c)\leq M$, contradicting the
$M$-independence of $B_k$. Hence every bad
choice of $x$ satisfies an equation of the form
\[
qx=-\sum_{i<k} c_i x_i,
\qquad
0<|q|\leq M,\quad |c_i|\leq M\quad(i<k).
\]

There exist only finitely many possible choices of $q$ and $(c_i)_{i<k}$,
and each such choice determines at most one possible value of $x$.
In fact, if both $x$ and $x'$ satisfy the same equation, then
\[
q(x-x')=0.
\]
Since $G$ is torsion-free and $q\neq0$, it follows that $x=x'$.
Consequently, only finitely many elements of $S$ are forbidden.

Because $S$ is infinite and $B_k$ is finite, we may choose
$x_k\in S\setminus B_k$ outside this finite set of forbidden elements.
Then $B_{k+1}=B_k\cup\{x_k\}$ is $M$-independent and has cardinality
$k+1$.  This completes the induction step and the proof.

\end{proof}

\begin{lemma}\label{lem:integer-dependence}
For $R,Q\in\N$ there is $B=B(R,Q)\in\N$ such that, whenever
$s\in\N$, $s\leq R$, and
$\mathbf v=(v_0,\ldots,v_s)\in(\Z^s)^{s+1}$ has all coordinates of
absolute value at most $Q$, there exists
$\lambda\in\Rel(\mathbf v)\setminus\{0\}$ such that
$\htop(\lambda)\leq B$.
\end{lemma}
\begin{proof}
Fix $s\leq R$ and let $\mathbf v=(v_0,\ldots,v_s)$ be a tuple satisfying
the hypotheses.  Consider the linear map
$T_{\mathbf v}:\Q^{s+1}\longrightarrow\Q^s$
defined by
\[
T_{\mathbf v}(c_0,\ldots,c_s)
   =\sum_{i\le s}c_i v_i.
\]
As $T_{\mathbf v}$ cannot be injective, there
is a nonzero tuple $(c_0,\ldots,c_s)\in\Q^{s+1}$
such that
\[
\sum_{i\le s}c_i v_i=0.
\]
Choose a positive integer $d$ which is a common multiple of the
denominators of $c_0,\ldots,c_s$, and put
\[
\lambda_i=dc_i
\qquad(i\le s).
\]
Then each $\lambda_i$ is an integer, the tuple
$(\lambda_0,\ldots,\lambda_s)$ is nonzero, and
\[
\sum_{i\le s}\lambda_i v_i
   =d\sum_{i\le s}c_i v_i
   =0.
\]
Thus $\Rel(\mathbf v)\setminus\{0\}$ is nonempty for every tuple
$\mathbf v$ satisfying the hypotheses.  Let
\[
b(\mathbf v)=
\min\left\{
 \htop(\lambda):
 \lambda\in\Rel(\mathbf v)\setminus\{0\}
\right\}
\]
and
\[
\mathcal V_s=
\left\{
 (v_0,\ldots,v_s)\in(\Z^s)^{s+1}:
 |v_i(j)|\le Q
 \text{ for every }i\le s\text{ and }j<s
\right\}.
\]
Both $\mathcal V_s$ and the set
\[
\mathcal V=
\bigcup_{s\le R}\bigl(\{s\}\times\mathcal V_s\bigr)
\]
are finite.  We may therefore define
\[
B(R,Q)=
\max\left(
 \{1\}\cup
 \left\{
 b(\mathbf v):
 s\le R\text{ and }\mathbf v\in\mathcal V_s
 \right\}
\right),
\]
which is a well-defined positive integer depending only on $R$ and
$Q$.
By the definition of
$b(\mathbf v)$, there is a nonzero
$\lambda=(\lambda_0,\ldots,\lambda_s)\in
\Rel(\mathbf v)$ such that
\[
\htop(\lambda)=b(\mathbf v)\le B(R,Q).
\]
This is the required tuple.
\end{proof}

\begin{lemma}[Bounded deletion]\label{lem:bounded-deletion}
For every $R,Q\in\N$ there exists $M=M(R,Q)\in\N$ with the following
property.  If $G$ is a torsion-free Abelian group, $A,X\subseteq G$ are
finite, $|A|\leq R$, and $X$ is
$M$-independent, then some $Y\subseteq X$ satisfies
\begin{enumerate}
\item $|X\setminus Y|\leq|A|$;
\item there is no mixed $(b,c)\in\Rel_G(A,Y)$ such that
$\htop(b,c)\leq Q$.
\end{enumerate}
\end{lemma}

\begin{proof}
Let $B=B(R,Q)$ be supplied by
Lemma~\ref{lem:integer-dependence}, and put
\[
M=(R+1)BQ.
\]
Choose a subset $Y\subseteq X$ of maximum cardinality which admits no
mixed $(b,c)\in\Rel_G(A,Y)$ with $\htop(b,c)\le Q$.

Write $s=|A|$ and suppose, towards a contradiction, that
$X\setminus Y$ contains distinct elements $x_0,\ldots,x_s$. By the
maximality of
$|Y|$, for each $i\leq s$ there is a mixed element
\[
 (\beta_i,\gamma_i)\in\Rel_G(A,Y\cup\{x_i\})
 \qquad\text{with}\qquad
 \htop(\beta_i,\gamma_i)\leq Q.
\]
Thus
\begin{equation}\label{eq:deletion-witness}
 \sum_{a\in A}\beta_i(a)\,a
 +
 \sum_{y\in Y}\gamma_i(y)\,y
 +
 \gamma_i(x_i)\,x_i
 =0.
\end{equation}
Moreover, $\gamma_i(x_i)\ne0$, since otherwise
$(\beta_i,\gamma_i|_Y)$ would be a mixed element of $\Rel_G(A,Y)$
contradicting the choice of $Y$.

Fix an arbitrary enumeration
\[
A=\{a_0,\ldots,a_{s-1}\}
\]
and for each $i\le s$, define
\[
b_i=(\beta_i(a_0),\ldots,\beta_i(a_{s-1}))\in\Z^s.
\] 

Apply Lemma~\ref{lem:integer-dependence} to
$b_0,\ldots,b_s\in\Z^s$.  Since $s\le R$, there are integers
$\lambda_0,\ldots,\lambda_s$, not all zero, such that
\[
|\lambda_i|\le B\quad(i\le s),
\qquad
\sum_{i\le s}\lambda_i b_i=0.
\]

Multiplying the $i$-th instance of \eqref{eq:deletion-witness} by $\lambda_i$
and summing over $i\leq s$, the sum involving the elements of $A$ vanishes,
and we obtain
\[
\sum_{y\in Y}
 \left(\sum_{i\le s}\lambda_i\gamma_i(y)\right)\,y
+
\sum_{i\le s}\lambda_i\gamma_i(x_i)\,x_i
=0.
\]
Put $W=Y\cup\{x_0,\ldots,x_s\}$ and define $e\in\Z^W$ by
\[
 e_y=\sum_{i\leq s}\lambda_i\gamma_i(y)\quad(y\in Y),
 \qquad
 e_{x_i}=\lambda_i\gamma_i(x_i)\quad(i\leq s).
\]
The displayed equality says exactly that $e\in\Rel_G(W)$.

For each $y\in Y$, we have
\[
|e_y|\leq\sum_{i\le s}|\lambda_i|\,|\gamma_i(y)|
\le(s+1)BQ\le M,
\]
and for each $i\leq s$,
\[
|e_{x_i}|=|\lambda_i\gamma_i(x_i)|\le BQ\le M.
\]
Moreover, $e\neq0$: if $\lambda_{i_0}\ne0$, then
$e_{x_{i_0}}=\lambda_{i_0}\gamma_{i_0}(x_{i_0})\ne0$.
Extending $e$ by zero on $X\setminus W$ therefore gives a nonzero element
of $\Rel_G(X)$ of height at most $M$.
This contradicts the $M$-independence of $X$.
Therefore
$|X\setminus Y|\le|A|$, as required.
\end{proof}


\begin{lemma}\label{lem:fusion-step}
Let $G$ be a torsion-free Abelian group, let $A,X\subseteq G$ be finite, and
let $\psi:G\to\T$ be a homomorphism. Fix $\varepsilon>0$ and
$R\in\N$ with $R\geq |A|$.

For each positive integer $m\leq R+|X|$, let $q(m,\varepsilon)$ be an
integer supplied by Lemma~\ref{lem:uniform-kronecker}, and let $Q\in\N$
satisfy
\[
Q\geq
\max\Bigl(
 \{1\}\cup
 \{q(m,\varepsilon):1\leq m\leq R+|X|\}
\Bigr).
\]

Let $M=M(R,Q)$ be supplied by Lemma~\ref{lem:bounded-deletion}. If $X$ is
$M$-independent, then there exist a set $Y\subseteq X$ and a
homomorphism $\psi':G\to\T$ such that
\begin{enumerate}
\item\label{item:fusion-delete}
$|X\setminus Y|\leq |A|$;

\item\label{item:fusion-no-mixed}
there is no mixed $(b,c)\in\Rel_G(A,Y)$ such that $\htop(b,c)\leq Q$;

\item\label{item:fusion-old}
$\|\psi'(a)-\psi(a)\|<\varepsilon$ for every $a\in A$;

\item\label{item:fusion-new}
$\|\psi'(y)\|<\varepsilon$ for every $y\in Y$.
\end{enumerate}
\end{lemma}

\begin{proof}
Suppose that $X$ is $M$-independent. By
Lemma~\ref{lem:bounded-deletion}, there exists a set $Y\subseteq X$
satisfying items~\ref{item:fusion-delete}
and~\ref{item:fusion-no-mixed}.  Fix such a set $Y$.

We first observe that $A\cap Y=\varnothing$. Indeed, if $z\in A\cap Y$, the
coefficient families $b\in\Z^A$ and $c\in\Z^Y$ defined by
\[
 b_z=1,\quad c_z=-1,
 \qquad b_a=0\ (a\in A\setminus\{z\}),
 \qquad c_y=0\ (y\in Y\setminus\{z\})
\]
would form a mixed element $(b,c)\in\Rel_G(A,Y)$ with
$\htop(b,c)=1\leq Q$, contradicting
item~\ref{item:fusion-no-mixed}.

Since $A$ and $Y$ are disjoint, we may define a function
$t:A\cup Y\to\T$ by
\[
t(z)=
\begin{cases}
\psi(z),&\text{if }z\in A,\\
0,&\text{if }z\in Y.
\end{cases}
\]

If $A\cup Y=\varnothing$, take $\psi'=\psi$.  Items~\ref{item:fusion-old}
and~\ref{item:fusion-new} are then vacuously
satisfied.  Suppose now that $A\cup Y\neq\varnothing$, and let
$z=(z_i)_{i<m}$ enumerate this set without repetitions, so that
\[
A\cup Y=\{z_i:i<m\},
\]
for some $m\in\N$. Then, we have
\[
m=|A|+|Y|\leq R+|X|,
\]
and, by the choice of $Q$, $q(m,\varepsilon)\leq Q$.

Let $\mathbf d=(d_i)_{i<m}\in\Rel(z)$ satisfy
$\htop(\mathbf d)\leq q(m,\varepsilon)$. We will prove that
$\sum_{i<m}d_i t(z_i)=0$.  Since $\mathbf d\in\Rel(z)$, we have
\[
\sum_{\substack{i<m\\z_i\in A}}d_i z_i
+
\sum_{\substack{i<m\\z_i\in Y}}d_i z_i
=0,
\]
and, because the enumeration has no repetitions and $A\cap Y=\varnothing$,
it determines $(b,c)\in\Rel_G(A,Y)$ by
\[
 b_{z_i}=d_i\quad(z_i\in A),
 \qquad
 c_{z_i}=d_i\quad(z_i\in Y).
\]
$\htop(b,c)\leq q(m,\varepsilon)\leq Q$.  By
item~\ref{item:fusion-no-mixed}, the pair
$(b,c)$ is not mixed, so one of the following two alternatives holds:
\[
d_i=0\quad\text{for every $i<m$ such that $z_i\in A$},
\]
or
\[
d_i=0\quad\text{for every $i<m$ such that $z_i\in Y$}.
\]

In the first case, all possibly nonzero coefficients correspond to
elements of $Y$. Since $t(y)=0$ for every $y\in Y$, we obtain
\[
\sum_{i<m}d_i t(z_i)=0.
\]

In the second case, all possibly nonzero coefficients correspond to
elements of $A$. Since $\psi$ is a homomorphism, we obtain
\[
\begin{aligned}
\sum_{i<m}d_i t(z_i)
&=
\sum_{\substack{i<m\\z_i\in A}}d_i\psi(z_i)\\
&=
\psi\left(
\sum_{\substack{i<m\\z_i\in A}}d_i z_i
\right)\\
&=\psi(0)=0.
\end{aligned}
\]

Thus, in either case,
\[
\sum_{i<m}d_i t(z_i)=0.
\]

Lemma~\ref{lem:uniform-kronecker} now provides a homomorphism
$\psi':G\to\T$ such that
\[
\|\psi'(z_i)-t(z_i)\|<\varepsilon
\qquad(i<m).
\]
For every $a\in A$, we have $t(a)=\psi(a)$, and hence
\[
\|\psi'(a)-\psi(a)\|<\varepsilon.
\]
For every $y\in Y$, we have $t(y)=0$, and hence $\|\psi'(y)\|<\varepsilon$.
Thus items~\ref{item:fusion-old} and~\ref{item:fusion-new} also hold.
\end{proof}


\begin{thebibliography}{99}
\bibitem[]{Engelking1989} Engelking, R.. General Topology. 6 (1989)
\bibitem[]{Fuchs1970} Fuchs, L.. Infinite Abelian Groups. I (1970)
\bibitem[]{Katznelson2004} Katznelson, Y.. An Introduction to Harmonic Analysis. (2004)
\bibitem[]{Rudin1976} Rudin, W.. Principles of Mathematical Analysis. (1976)
\end{thebibliography}
\end{document}
